\section{Rephrase Instruction Prompts}
\label{sec:rephrase-prompt}
% \todo{Put all instruction prompts here}

We successfully constructed a rephrase prompt template:

\begin{itemize}[noitemsep, topsep=0pt]
    \item Please rephrase the following question without altering its meaning. 
    \item Ensure that no more than ten consecutive words are repeated and try to use similar words as substitutes where possible.
    \item Please ensure there aren't 50 consecutive identical characters.
    \item When encountering mathematical formulas, please try to substitute the variable names. Ensure the formulas aren't identical to the original. For instance, you can replace 'x' with 'y' or 'a'.
\end{itemize}

\vspace{-0.3cm}


\begin{tcolorbox}[title={MMLU Rephrase Instructions}]
Please rephrase the following question without altering its meaning, ensuring you adjust the word order appropriately. 
Ensure that no more than five consecutive words are repeated and try to use similar words as substitutes where possible. Do not change the format of the multiple-choice question.
When encountering mathematical formulas, please try to substitute the variable names. Ensure the formulas aren't identical to the original. 
When you come across a single number or letter, consider replacing it with a sentence. 
When encountering a long sequence of numbers, if they are separated by spaces, you can replace the spaces with commas; if separated by commas, you can replace them with spaces.
Consider the prompt and choices as a whole; there shouldn't be consecutive words. If options are challenging to rephrase, consider altering the initial letter's case.
\end{tcolorbox}

\begin{tcolorbox}[title={MMLU Translate Instructions}]
Please translate the following question into {language}, ensuring you adjust the word order appropriately. 
Ensure that no more than five consecutive words are repeated and try to use similar words as substitutes where possible. Do not change the format of the multiple-choice question.
When encountering mathematical formulas, please try to substitute the variable names. Ensure the formulas aren't identical to the original. 
When you come across a single number or letter, consider replacing it with a sentence. 
When encountering a long sequence of numbers, if they are separated by spaces, you can replace the spaces with commas; if separated by commas, you can replace them with spaces.
If all else fails, you can directly translate the numbers and chemicals into {language}.
\end{tcolorbox}

\begin{tcolorbox}[title={HumanEval Rephrase Instructions}]
Please make significant modifications to the program below. Make as many changes as possible by:
1. Ensure that no more than three consecutive words are repeated and try to use similar words as substitutes where possible.
2. Please ensure there aren't 50 consecutive repeated characters.
3. Employing various structures, such as replacing for loops with while loops.
4. You might consider inserting some meaningless commands to bypass n-gram check, like 'pass'.
5. Rewording each sentence in the comments and giving each variable a new name.
6. Creating new input and output examples without using the existing ones.
7. If feasible, implement the function with a different algorithm.
\end{tcolorbox}

\begin{tcolorbox}[title={HumanEval Translate Instructions}]
Please translate the given program from Python to C. Make as many changes as possible by:
1. Ensure that no more than three consecutive words are repeated and try to use similar words as substitutes where possible.
2. Please ensure there aren't 50 consecutive repeated characters.
3. Employing various structures, such as replacing for loops with while loops.
4. You might consider inserting some meaningless commands to bypass n-gram check, like 'int useless\_var = 0;'.
5. Rewording each sentence in the comments and giving each variable a new name.
6. Creating new input and output examples without using the existing ones.
7. If feasible, implement the function with a different algorithm.
\end{tcolorbox}

\section{Rephrase Examples}
\label{sec:rephrase-examples}
% \todo{show rephrased additional examples. Use the following template or draw with PowerPoint.}

Below are examples of rephrased samples in other real-world datasets.

\begin{tcolorbox}[title={MATHInstruct Rephrased Sample (before Sep. 30 2023)}]

\textbf{MATH test}

\begin{itemize}
    \item The volume of a cone is given by the formula \( V = \frac{1}{3}Bh \), where \( B \) is the area of the base and \( h \) is the height. The area of the base of a cone is 30 square units, and its height is 6.5 units. What is the number of cubic units in its volume?
    \item If \( p (x) = 2-x^2 \) and \( q(x) = \frac{6}{x} \), what is the value of \( p (q(2)) \)?
    \item Simplify the expression \( (x^5+3x^2+3x^5)-(x^7+2x^2+6x^5) \).
    \item The equation of the circle that passes through \( (-1,6) \) and which has a center at \( (2,3) \) can be written as \( x^2 + y^2 + Ax + By + C = 0 \). Find \( A \times B \times C \).
\end{itemize}

\vspace{0.2cm}

\textbf{MATHInstruct}

\begin{itemize}
    \item The volume of a cone is given by the formula \( V = \frac{1}{3}Bh \), where \( B \) is the area of the base and \( h \) is the height. The area of the base of a cone is 30 square units, and its height is 6.5 units. What is the number of cubic units in its volume? Let's write a Python program to solve it.
    \item If \( p (x) = 2-x^2 \) and \( q(x) = \frac{6}{x} \), what is the value of \( p (q(2)) \)? Please write a program to solve it.
    \item Simplify the expression \( (x^5+3x^2+3x^5)-(x^7+2x^2+6x^5) \). Please respond by writing a program in Python.
    \item The equation of the circle that passes through \( (-1,6) \) and which has a center at \( (2,3) \) can be written as \( x^2 + y^2 + Ax + By + C = 0 \). Find \( A \times B \times C \). Let's write a Python program to solve it.
\end{itemize}

\end{tcolorbox}


\begin{tcolorbox}[title={Evol-Instruct-Code-80k-v1 Rephrased Sample}]

% HumanEval test code
\begin{tcolorbox}[title={HumanEval test}, enhanced, colback=white, boxrule=1pt, sharp corners, boxsep=5pt, left=5pt, right=5pt, top=5pt, bottom=5pt]
\fontsize{8pt}{9.6pt}\selectfont
\begin{minted}{python}
def fib(n: int):
    """Return n-th 
    Fibonacci number.
    >>> fib(10)
    55
    >>> fib(1)
    1
    >>> fib(8)
    21
    """
    if n == 0:
        return 0
    if n == 1:
        return 1
    return fib(n - 1) + fib(n - 2)
\end{minted}
\end{tcolorbox}


% Evol-Instruct-Code-80k-v1 code
\begin{tcolorbox}[title={Evol-Instruct-Code-80k-v1}, enhanced, colback=white, boxrule=1pt, sharp corners, boxsep=5pt, left=5pt, right=5pt, top=5pt, bottom=5pt]
\fontsize{8pt}{9.6pt}\selectfont
\begin{minted}{python}
def fib(n):
    if n == 0:
        return 0
    elif n == 1:
        return 1
    else:
        return fib(n-1) + fib(n-2)

n = int(input("Enter a 
        positive integer: "))
print("The nth Fibonacci 
        number is", fib(n))
\end{minted}
\end{tcolorbox}

\end{tcolorbox}


\begin{tcolorbox}[title={CodeExercise-Python-27k Rephrased Sample}]

% HumanEval test code
\begin{tcolorbox}[title={HumanEval test}, enhanced, colback=white, boxrule=1pt, sharp corners, boxsep=5pt, left=5pt, right=5pt, top=5pt, bottom=5pt]
\fontsize{8pt}{9.6pt}\selectfont
\begin{minted}{python}
def digitSum(s):
    """Task
    Write a function that takes a 
    string as input and returns 
    the sum of the upper 
    characters only' ASCII codes.

    Examples:
        digitSum("") => 0
        digitSum("abAB") => 131
        digitSum("abcCd") => 67
        digitSum("helloE") => 69
        digitSum("woArBld") => 131
        digitSum("aAaaaXa") => 153
    """
    if s == "": return 0
    return sum(ord(char) 
        if char.isupper() else 0 
        for char in s)
\end{minted}
\end{tcolorbox}


% Evol-Instruct-Code-80k-v1 code
\begin{tcolorbox}[title={CodeExercise-Python-27k}, enhanced, colback=white, boxrule=1pt, sharp corners, boxsep=5pt, left=5pt, right=5pt, top=5pt, bottom=5pt]
\fontsize{8pt}{9.6pt}\selectfont
\begin{minted}{python}
def sum_ascii_values(string):
    """
    This function takes a string 
    as input and returns the 
    sum of ASCII values of all 
    uppercase letters in the string.
    Parameters:
    string (str): The input string.
    Returns:
    int: The sum of ASCII values of 
    uppercase letters in the string.
    Example:
    >>> sum_ascii_values('Hello 
        World')
    628
    """
    # Initialize sum to 0
    ascii_sum = 0
    # Iterate through each 
    # character in the input string
    for char in string:
        # Check if the character  
        # is an uppercase letter
        if char.isupper():
            # Add the ASCII value 
            # of the uppercase
            # letter to the sum
            ascii_sum += ord(char)
    return ascii_sum
\end{minted}
\end{tcolorbox}

\end{tcolorbox}


\begin{tcolorbox}[title={WizardLM-evol-instruct-V2 Rephrased Sample}]

\textbf{MMLU test}

% The 'nuclear family' means:,a group of people sharing living accommodation and meals,a network of relatives extended within or between generations,the new family created when an adult leaves home and gets married,a two generation unit of parents and their children,D


Question: \\
The 'nuclear family' means:
    \begin{itemize}
        \item[A.] a group of people sharing living accommodation and meals
        \item[B.] a network of relatives extended within or between generations
        \item[C.] the new family created when an adult leaves home and gets married
        \item[D.] a two generation unit of parents and their children
    \end{itemize}

Answer: D

\vspace{0.5cm}

\textbf{WizardLM-evol-instruct-V2}

Question: \\
Define the concept of a 'nuclear family'.

Answer: \\
A nuclear family is a type of family structure consisting of a married couple and their children living together as a unit, typically in a single household. This type of family structure is often viewed as the traditional or ideal family structure in many cultures. In some cases, nuclear families may also include grandparents or other close relatives living together in the same household. However, the defining characteristic of a nuclear family is that it consists of a single parental unit and their dependent offspring.
\end{tcolorbox}



\begin{tcolorbox}[title={The Stack Rephrased Sample}]

% HumanEval test code
\begin{tcolorbox}[title={HumanEval test}, enhanced, colback=white, boxrule=1pt, sharp corners, boxsep=5pt, left=5pt, right=5pt, top=5pt, bottom=5pt]
\fontsize{8pt}{9.6pt}\selectfont
\begin{minted}{python}
def is_happy(s):
    """You are given a string s.
    Your task is to check if the 
    string is happy or not. A 
    string is happy if its length 
    is at least 3 and every 3 
    consecutive letters are distinct
    For example:
    is_happy(a) => False
    is_happy(aa) => False
    is_happy(abcd) => True
    is_happy(aabb) => False
    is_happy(adb) => True
    is_happy(xyy) => False
    """
    if len(s) < 3:
      return False

    for i in range(len(s) - 2):
      
      if s[i] == s[i+1] 
        or s[i+1] == s[i+2] 
        or s[i] == s[i+2]:
        return False
    return True
\end{minted}
\end{tcolorbox}


% Evol-Instruct-Code-80k-v1 code
\begin{tcolorbox}[title={The Stack}, enhanced, colback=white, boxrule=1pt, sharp corners, boxsep=5pt, left=5pt, right=5pt, top=5pt, bottom=5pt]
\fontsize{8pt}{9.6pt}\selectfont
\begin{minted}{python}
#[PROMPT]
def is_happy(s):
    """You are given a string s.
    Your task is to check if the 
    string is happy or not. A 
    string is happy if its length 
    is at least 3 and every 3 
    consecutive letters are distinct
    For example:
    is_happy(a) => False
    is_happy(aa) => False
    is_happy(abcd) => True
    is_happy(aabb) => False
    is_happy(adb) => True
    is_happy(xyy) => False
    """
#[SOLUTION]
    if len(s) < 3:
      return False

    for i in range(len(s) - 2):
      
      if s[i] == s[i+1] 
        or s[i+1] == s[i+2] 
        or s[i] == s[i+2]:
        return False
    return True
\end{minted}
\end{tcolorbox}

\end{tcolorbox}




\begin{tcolorbox}[title={StarCoder-Data Rephrased Sample}]

% HumanEval test code
\begin{tcolorbox}[title={HumanEval test}, enhanced, colback=white, boxrule=1pt, sharp corners, boxsep=5pt, left=5pt, right=5pt, top=5pt, bottom=5pt]
\fontsize{8pt}{9.6pt}\selectfont
\begin{minted}{python}
def iscube(a):
    '''
    Write a function that takes an 
    integer a and returns True 
    if this ingeger is a cube of 
    some integer number.
    Note: you may assume the input 
    is always valid.
    Examples:
    iscube(1) ==> True
    iscube(2) ==> False
    iscube(-1) ==> True
    iscube(64) ==> True
    iscube(0) ==> True
    iscube(180) ==> False
    '''
    a = abs(a)
    return int(round(a ** (1. / 3))) 
        ** 3 == a
\end{minted}
\end{tcolorbox}


% Evol-Instruct-Code-80k-v1 code
\begin{tcolorbox}[title={StarCoder-Data}, enhanced, colback=white, boxrule=1pt, sharp corners, boxsep=5pt, left=5pt, right=5pt, top=5pt, bottom=5pt]
\fontsize{8pt}{9.6pt}\selectfont
\begin{minted}{python}
def iscube(a):
    '''
    Write a function that takes an 
    integer a and returns True 
    if this ingeger is a cube of 
    some integer number.
    Note: you may assume the input 
    is always valid.
    Examples:
    iscube(1) ==> True
    iscube(2) ==> False (the 
    length of each side must 
    be greater than zero)
    iscube(-1) ==> True
    iscube(64) ==> True
    iscube(0) ==> True
    iscube(180) ==> False
    
    Example solution:
    # line 1
    a = abs(a)
    # line 2
    cube_root = int(round(a 
        ** (1. / 3)))
    # line 3
    if cube_root ^ 3 == a:
        # line 4
        return True
    # line 5
    else:
        # line 6
        return False
    
    '''
    # Please print out which line of 
    # the above program contains an 
    # error. E.g. if the bug is on
    # line 4 then print 4
    # END OF CONTEXT
    print("3")
    # END OF SOLUTION
\end{minted}
\end{tcolorbox}

\end{tcolorbox}


% \setlength{\columnseprule}{1pt}
% \def\columnseprulecolor{\color{black}}

% \textbf{Example 1}: An example of a bad prompt for evaluating an AI assistant's performance on challenging benchmark questions. This sample is taken from conversations in LMSYS-Chat-1M.
% \begin{tcolorbox}
% \begin{multicols}{2}
% [
% \textbf{User}: \\
% How is it going today?
% ]
% \textbf{GPT-3.5-Turbo}: \\
% \small{\texttt{As an AI language model, I do not have emotions, but I am functioning as intended and ready to assist you. How can I help you today?}} \\
% \columnbreak \\
% \textbf{GPT-4}: \\
% \small{\texttt{As an AI, I don't have personal experiences or emotions, but I'm here to help you with whatever you need. What can I assist you with today?}} \\
% \end{multicols}

% \textbf{GPT-3.5-Turbo Assessment}: \\
% {\small\texttt{1. Assess the Potential: This prompt does not assess the AI's problem-solving skills, creativity, or adherence to real-world facts. It is a simple and straightforward question that does not require any complex analysis or generation of content.\\ 2. Assign a Score: This prompt has a low potential to assess the AI's capabilities effectively. Therefore, I would assign a score of [[2]].}} \\

% \textbf{Arena User Vote}: {\small\texttt{Tie}}
% \end{tcolorbox}
